\chapter{Conclusão}
O presente trabalho demonstrou, de forma concreta e bem-sucedida, a viabilidade do desenvolvimento de um sistema embarcado de reconhecimento facial em tempo real em \textit{\textit{hardware}} reconfigurável. O projeto integrou conhecimentos de aprendizado de máquina, arquitetura de computadores digitais, comunicação serial e processamento de vídeo em uma única solução de engenharia, atuando de ponta a ponta.

A opção inicial pela rede binária revelou-se inviável diante da necessidade de identificar individualmente cada membro autorizado e exibir seu nome no monitor VGA. A adoção de uma arquitetura \textit{Tiny-CNN} com 19 classes — 18 membros autorizados e uma classe dedicada ao reconhecimento de desconhecidos — resolveu essa limitação. Além disso, permitiu tratar a rejeição de indivíduos não cadastrados como uma classe real e treinável, em vez de recorrer a limiares arbitrários de confiança externos. Essa decisão arquitetural refletiu uma compreensão adequada do problema, pois qualquer indivíduo não cadastrado precisa ser corretamente identificado como tal antes mesmo de se tentar qualquer autorização de acesso.

A etapa de construção do \textit{dataset} e de treinamento do modelo evidenciou a importância do equilíbrio e da diversidade dos dados de entrada. As técnicas de aumento de dados (\textit{data augmentation}) foram essenciais para compensar a limitação natural do volume de capturas e aumentar a capacidade de generalização do algoritmo. A classe de desconhecidos recebeu atenção especial, com a integração de conjuntos de dados externos e imagens sintéticas, o que foi fundamental para minimizar falsos positivos e cobrir a enorme variabilidade do rosto humano fora do escopo de membros cadastrados. O esquema de ponderação diferenciada da função de perda refletiu com precisão a assimetria de custo do problema real: negar acesso a um membro cadastrado é, operacionalmente, mais grave do que classificar um desconhecido como tal.

A quantização dos pesos do modelo treinado para o formato de ponto fixo representou o elo crítico entre o domínio de \textit{software} e o de \textit{hardware}. A estratégia de regularização adotada durante o treinamento cumpriu um papel duplo: além de mitigar o \textit{overfitting}, contribuiu diretamente para a redução do erro de quantização posterior. Isso evidenciou a conexão entre as decisões de treinamento em \textit{software} e as restrições físicas do dispositivo — destacando-se como uma das escolhas mais relevantes do projeto.

A arquitetura de \textit{hardware}, implementada em Verilog, refletiu um domínio consistente sobre os recursos físicos do dispositivo e suas limitações. A máquina de estados finitos orquestrou da forma esperada todos os módulos do sistema, garantindo a correta sequência de operações desde o recebimento da imagem até a saída da inferência.

A adoção do coeficiente de Pearson como métrica principal revelou uma compreensão precisa do que realmente importa para a fidelidade de uma rede neural embarcada: não é necessário que os valores absolutos das ativações intermediárias sejam idênticos entre \textit{software} e \textit{hardware}, mas sim que a proporcionalidade relativa entre eles seja mantida. Os resultados obtidos confirmaram essa premissa com solidez, mantendo uma correlação superior a 0,9997 em todas as camadas analisadas. 

A integração do protocolo de comunicação serial como canal de alimentação de dados dinâmicos foi feita de forma unilateral, tendo apenas o FPGA como receptor dessas informações. O protocolo de controle de fluxo estabeleceu uma saída satisfatória para o roteamento dos dados entre os dois modos de operação do sistema: exibição de vídeo contínuo e inferência facial. O subsistema de exibição em monitor completou o ciclo de \textit{feedback} visual de forma eficaz, compondo em tempo real a imagem capturada com os \textit{sprites} de identificação dos membros reconhecidos.

No tocante às dificuldades enfrentadas, a discrepância no volume de imagens entre os integrantes da equipe evidenciou um desafio prático comum em projetos de reconhecimento facial com dados reais: a quantidade e a qualidade dos dados de treinamento por indivíduo impactam diretamente a confiabilidade da classificação para aquela classe específica, mesmo quando técnicas de aumento de dados são empregadas. Esse desequilíbrio, embora não tenha comprometido o funcionamento geral do sistema, serve como lição sobre a importância de padronizar os critérios de coleta desde o início.

Em síntese, o projeto integrou com êxito os quatro artefatos propostos no cronograma --- aquisição e treinamento \textit{offline}, arquitetura RTL, \textit{pipeline} de visão e comunicação, e integração final com demonstração prática ---, resultando em um sistema coerente e validado quantitativamente. A Fechadura Biométrica Inteligente desenvolvida neste trabalho prova que é possível, com critérios rigorosos de projeto e uma metodologia gradual e incremental, realizar a inferência de redes neurais convolucionais em \textit{hardware} reconfigurável de baixo custo. Isso abre caminho para aplicações de inteligência artificial embarcada em sistemas de segurança autônomos e de baixa latência.
